\documentclass{jlreq}
\begin{document}

与えられたグラフに対して、任意の2頂点間の最短経路を計算するアルゴリズムを書いた。

入力：グラフを表す隣接頂点関数T、頂点${v,w}$間の枝の重みを返す関数$l_{v,w}$、始点$s$、終点$t$

出力：最短経路を表す頂点列

最短経路探索のアルゴリズムを以下に示す。

$s$から頂点$v$までの最短経路長を表す関数$d_v$

$s$から頂点$v$までの最短経路を考えたとき、$v$の直前の頂点を表す関数$prev_v$

展開する頂点を格納する集合$A$

を用意する。

1. $d_s=0$, $s$以外のすべての頂点に対し$d_v=∞$とする。

2. $s$を$A$に追加。

3. $A$に含まれる頂点のうち$d_v$の値が最小となる頂点$v$を取り出す。（$A$からは消す）

4. $v$に隣接する各頂点$v$に対して、

if $d_w=∞$ then $d_w=d_v+l_{v,w}$とし$w$を$A$に追加。$prev_w=v$
elseif $d_w > d_v+l_{v,w}$ then $d_w=d_v+l_{v,w}$, $prev_w=v$

5. $A$が空なら手順6へ、空でないなら手順3へ戻る。
6. $prev_t$から辿って行けば$t$から$s$への経路が得られるのでこれを出力する。最短経路長は$d_t$で与えられる。

これで最短距離が求まる正当性を簡単に説明する。

まず各頂点は「距離が確定した頂点」と「距離が確定していない頂点」に分類できる。
「距離が確定した頂点」は手順3で$A$から取り出されたことのある頂点集合である。
手順3を行う直前の時点ですべての$d_v$は、「距離が確定した頂点」から直接その頂点に移動したときの最短距離以下になっている。（手順4により）
実はここで「距離が確定していない頂点」のうち、現時点でもっとも$d_v$の値が小さい頂点は、その$d_v$の値を確定させて良い。
手順4のような更新を喰らうことがもうないため、これ以上$d_v$を小さくすることができないからである。
このようにして最短距離がドミノ式に確定していく。

$A$はヒープによって実装した。ヒープの各頂点は「頂点番号」と「最短距離」という２次元の情報を持ち、「最短距離」によって比較される。
$A$の内部にある最短距離情報を直接更新する必要はなく、新たな「頂点番号」「最短距離」の２次元の値をヒープに追加すれば上記のアルゴリズムが正当に動く。

\end{document}
